% Part: turing-machines
% Chapter: machines-computations
% Section: variants

\documentclass[../../../include/open-logic-section]{subfiles}

\begin{document}

\olfileid{tur}{mac}{var}
\olsection{ट्यूरिंग यंत्रांचे विविध प्रकार}

खरे तर ट्यूरिंग यंत्रांची व्याख्या अनेक प्रकारे करता येते;
आपली व्याख्या हा त्यांपैकी एकच प्रकार आहे. काही बाबतींत ती
इतर व्याख्यांपेक्षा अधिक मोकळी आहे. आपण कोणतीही सांत वर्णमाला
चालू देतो; अधिक मर्यादित व्याख्येत फितीवरील फक्त दोन चिन्हे,
$\TMstroke$ आणि~$\TMblank$, चालतील. आपण यंत्राला एकाच वेळी
फितीवर चिन्ह लिहू आणि पुढे सरकू देतो; इतर व्याख्यांत लिहिणे
किंवा सरकणे यांपैकी एकच क्रिया करता येते. शीर्ष न सरकवता
लिहिण्याची शक्यताही आपण ठेवतो; इतर व्याख्यांत $\TMstay$
ही ``सूचना'' नसते. दुसरीकडे, काही बाबतींत आपली व्याख्या अधिक
मर्यादित आहे. फीत फक्त एका दिशेला अनंत आहे असे आपण मानले;
इतर व्याख्यांत ती डावीकडे आणि उजवीकडे दोन्ही दिशांना अनंत
असते. एवढेच नव्हे, तर कितीही स्वतंत्र फिती किंवा घरांची
अनंत जाळीही चालू देता येईल. आपण ट्यूरिंग यंत्राचा सूचनासंच
संक्रमण फलनाने दर्शवतो; इतर व्याख्यांत संक्रमण संबंध वापरतात,
ज्यामुळे एखाद्या परिस्थितीत यंत्रासाठी एकाहून अधिक सूचना
शक्य असतात.

व्याख्येतील हा शेवटचा बदल विशेष लक्षवेधक आहे. आपल्या व्याख्येत
यंत्र~$q$ अवस्थेत~$\sigma$ चिन्ह वाचत असताना
$\delta(q, \sigma)$ पुढील चिन्ह, अवस्था आणि शीर्षाची जागा
निश्चित करते. पण सूचनासंच हा सध्याच्या अवस्था-चिन्ह जोड्या
$\tuple{q,
  \sigma}$ आणि पुढील अवस्था-चिन्ह-दिशा त्रयी $\tuple{q', \sigma',
  D}$ यांमधील संबंध म्हणून घेतला, तर यंत्राची कृती एकमेवपणे
निश्चित होणार नाही: सूचना-संबंधात $\tuple{q,
  \sigma, q', \sigma', D}$ आणि $\tuple{q, \sigma, q'', \sigma'', D'}$
या दोन्ही नोंदी असू शकतात. अशा वेळी आपल्याला
\emph{अनिर्धारक} ट्यूरिंग यंत्र मिळते. संगणकीय गुंतागुंत
सिद्धांतात अशा यंत्रांना महत्त्वाचे स्थान आहे.

ट्यूरिंग यंत्र कधी थांबते याबाबतही वेगवेगळे संकेत आहेत:
संक्रमण फलन अपरिभाषित असेल तेव्हा ते थांबते असे आपण मानतो;
इतर व्याख्यांत यंत्राने खास निर्दिष्ट थांबण्याच्या अवस्थेत
असणे आवश्यक असते. अशी निर्दिष्ट अवस्था मागितल्याने ट्यूरिंग
यंत्रांना काय संगणित करता येते यावर निर्बंध येत नाही, हे आपण
\olref[hal]{sec} मध्ये स्पष्ट केले आहे. आपल्या यंत्रांची फीत
एकाच दिशेला अनंत असल्यामुळे कधी एखादी सूचना नीट अमलात
आणता येत नाही: यंत्र सर्वांत डावीकडील घर वाचत आहे आणि
डावीकडे सरकायची सूचना आहे. आपल्या व्याख्येनुसार शीर्ष
फितीबाहेर ``पडण्याऐवजी'' त्याच घरावर राहते; पण त्या वेळी
यंत्र थांबेल अशीही व्याख्या करता आली असती. ती व्याख्याही
सममूल्य आहे: घर~$0$ वरून डावीकडे जाण्याचा प्रयत्न केल्यावर
थांबणाऱ्या यंत्राचे वर्तन अनुकरण करण्यासाठी प्रत्येक
$\delta(q,\TMendtape) =
\tuple{q',\sigma,\TMleft}$ संक्रमण काढून टाकता येईल.
मग~$\TMendtape$ वरून डावीकडे जाण्याचा प्रयत्न करण्याऐवजी
यंत्र थांबेल.\footnote{ही पद्धत \emph{अगदी तशी} चालत नाही,
कारण घर~$0$ सोडून इतर घरांतही $\TMendtape$ लिहिण्यास किंवा
वाचण्यास मनाई नाही (\olref[una]{ex:mover} पाहा). अशा वापरासाठी
दुसरे~$\TMendtape'$ चिन्ह जोडून ही अडचण दूर करता येते.}

संख्या निरूपित करण्याचेही वेगवेगळे मार्ग आहेत (आणि त्यामुळे
ट्यूरिंग यंत्र संगणित करत असलेल्या आदान-प्रदान फलनाचे रूपही
बदलते): आपण एकचिन्ही निरूपण वापरतो, पण द्विमान निरूपणही
वापरता येते. त्यासाठी $\TMblank$ आणि~$\TMendtape$ यांच्या
जोडीला आणखी दोन चिन्हे लागतात.

आता एक महत्त्वाची गोष्ट: यांपैकी कोणत्याही बदलामुळे नेमकी
कोणती फलने ट्यूरिंग-संगणनक्षम आहेत हे बदलत नाही.
\emph{एका व्याख्येनुसार एखादे फलन ट्यूरिंग-संगणनक्षम असेल,
तर त्या सर्व व्याख्यांनुसार ते ट्यूरिंग-संगणनक्षम असते.}

याची सविस्तर पडताळणी येथे करणार नाही. एक उदाहरण पाहू:
दोन दिशांना अनंत असलेली फीत किंवा अनेक फिती वापरू दिल्याने
संगणनशक्ती वाढत नाही. कारण साधारणपणे एकाच दिशेला अनंत असलेली
एक फीत वापरणारे ट्यूरिंग यंत्र अनेक फितींचे किंवा दोन्ही
दिशांना अनंत असलेल्या फितींचे अनुकरण करू शकते. उदाहरणार्थ,
अतिरिक्त अवस्था आणि सूचना वापरून अनेक फितींच्या किंवा
दोन दिशांना अनंत असलेल्या फितीच्या यंत्राचा कार्यक्रम, एका
दिशेला अनंत असलेली एकच फीत वापरणाऱ्या यंत्राच्या कार्यक्रमात
``रूपांतरित'' करता येतो. रूपांतरित यंत्र सम क्रमांकांची घरे
फीत~$1$ वरील घरांसाठी (किंवा द्विदिश अनंत फितीच्या ``धन''
बाजूच्या घरांसाठी), आणि विषम क्रमांकांची घरे फीत~$2$ वरील
घरांसाठी (किंवा त्या फितीच्या ``ऋण'' बाजूच्या घरांसाठी)
वापरू शकते.

\end{document}
